\section{Software architecture and execution model}\label{sec:architecture}

Figure~\ref{fig:architecture} summarizes the central design.  A
\texttt{MethodSpec} and its analytic coefficient algebra are interpreted by one
EGLM engine.  Numerical $\varphi$ and exponential actions are delegated through
a backend boundary.  Startup, adaptivity, backend selection and instrumentation
are orthogonal services around that path; they are not embedded in 47 separate
method implementations.

\begin{figure}[t]
  \centering
  \resizebox{\linewidth}{!}{\input{fig01_architecture_inline.tex}}
  \caption{Release architecture of \pkg{}.  Mathematical method definitions
  enter a common operator-coefficient algebra and explicit EGLM engine.  The
  engine requests actions from interchangeable numerical backends, while
  startup, adaptive control, dispatch and statistics remain separate layers.}
  \Description{Block diagram in which MethodSpec feeds an operator-coefficient
  algebra and a common EGLM engine.  The engine connects to dense, diagonal,
  Arnoldi, KIOPS-type, Leja and hybrid action backends.  Startup, adaptive
  control, automatic backend selection and statistics are shown as surrounding
  services rather than method-specific branches.}
  \label{fig:architecture}
\end{figure}

\subsection{Problem, state and solution objects}

A \texttt{SemilinearProblem} contains $L$, $N(t,y)$, the initial vector, a time
interval and, optionally, an exact solution used for validation or exact
startup.  Its minimal contract intentionally does not impose a PDE mesh,
Jacobian representation or spectral decomposition.  An \texttt{ExternalState}
stores the time and the $r\times n$ external-state array required by the chosen
GLM; its first row is always the physical state.  A \texttt{Solution} stores the
sampled times and states together with the method and backend identifiers,
step metadata and a statistics mapping.

This separation prevents a common source of hidden coupling.  A method does
not own an $L$ representation, and a backend does not own the nonlinear model.
The only place where both meet is the action context created for a particular
step size.

\subsection{The common step function}

The implementation of \texttt{step\_external} follows
Eqs.~\eqref{eq:eglm-stage}--\eqref{eq:eglm-output} directly.  At each stage it
first forms the action of the $U$ row on the external state and, for later
stages, the action of the strict-lower $A$ row on previously computed nonlinear
vectors.  It evaluates $N$ once at the stage node, stores $hN$, and then forms
the external outputs from $V$ and $B$.

The action-sum helper first asks the bound backend for an optional native
fused-sum operation.  If that operation is not available, it falls back to
individual analytic-expression actions.  Thus backend-specific fusion is
optional without changing method semantics.  This same mechanism is responsible
for the controlled fusion ablation in Section~\ref{sec:fusionexperiment}, where
the unfused path of a backend is obtained by hiding the fused operation from
its bound context.

\subsection{History and startup}

A one-step method has $r=1$ and needs no history.  For a method whose
external layout stores nonlinear history, $r-1$ previous nonlinear values must
be available before the main method can advance.  \pkg{} makes that
initialization policy explicit.  \texttt{ExactStartup} uses the problem's exact
solution and is intended primarily for validation.  \texttt{OneStepStartup}
uses any one-step \texttt{MethodSpec}; HochOst4 is the default numerical starter
when no exact solution is supplied.

The numerical starter bounds the accuracy a history method can reach in
general.  Each startup interval is advanced with a single step of the starter,
so the starting values carry its local error $O(h^{q+1})$, $q$ the order of the
starter.  Under the usual regularity and stability assumptions, and for a fixed
spatial discretization, the end-to-end error of a method of order $p$ is then
bounded by $C_1h^p+C_2h^{q+1}$.  This provides an order guarantee
$\min(p,q+1)$, hence $\min(p,5)$ with HochOst4; it does not impose a
universal upper bound on an observed order when leading terms cancel.  Section~\ref{sec:startup} measures this for the order six and seven
entries of the catalogue.  Making startup visible therefore matters both for accuracy
and for work comparisons: a benchmark that uses exact history isolates the
asymptotic method body, and it must not present that cost, or that order, as a
complete practical solve.  This distinction is stated explicitly for the
fixed-window work--precision experiment in Section~\ref{sec:methodwp}.

\subsection{Work accounting}

A backend that supplies a matrix-vector counter reports its cumulative
value in the top-level \texttt{matvecs} entry of its statistics.  Selectors that delegate to
several engines report the sum there and keep per-engine breakdowns in nested
entries, which the adaptive solvers never add to the total.  The work of an
attempt is the increment of that counter; when a backend reports no counter,
for example an exact diagonalization, the work is recorded as missing rather
than replaced by another quantity.  Nonlinear evaluations are counted by
wrapping the user's function, so that failed attempts are included.  These
rules are what make the work observables of Section~\ref{sec:experiments}
comparable across backends.

\subsection{Interface usage fragment}\label{sec:example}

The following fragment illustrates the solver interface once a semilinear
\texttt{problem} and the listed PyEXPINT objects have been imported or
constructed.  It is not a standalone program: the spatial grid, initial state
and imports belong to the complete example in the software companion.  No
exact solution is supplied for this use of numerical startup.
\begin{verbatim}
reference = solve_fixed(
    problem, HOCHOST4, 2/512,
    KiopsBackend(tol=1e-12)).y[-1]
multistep = solve_fixed(
    problem, ABNORSETT4, 2/128,
    KiopsBackend(tol=1e-11))
adaptive = solve_etd34(
    problem, KiopsBackend(tol=1e-11),
    options=ETD34Options(rtol=1e-8, atol=1e-10))
print(adaptive.stats)
\end{verbatim}
The reference here is a numerical comparison, not an exact solution or a
certified error bound.  The complete program is supplied as
\path{reproducibility/no_exact_history.py}, byte-identical to
\path{software/examples/no_exact_history.py} in the private companion.
Its clean-installation log records a comparison difference of
$6.53\times10^{-9}$ for ABNorsett4 and $7.39\times10^{-9}$ for ETD34,
with 56 accepted and one rejected adaptive step, in the recorded environment.
These are example-validation outputs, not a performance benchmark.

A method is added by constructing a \texttt{MethodSpec} from its nodes and
analytic coefficient arrays.  A backend provides \texttt{bind(L,h)}, returning
an action context, and \texttt{stats()}, with an authoritative top-level
\texttt{matvecs} total when that metric is available.  The adaptive interface
also requests \texttt{with\_tolerance(tol)}.  A callable hook takes precedence over the built-in cloning cases and is
responsible for preserving the required non-tolerance configuration.
Independent mutable statistics are needed when a new stateful work backend
is returned; the supplied protocol also permits a stateless realization to
return itself.  Each solve resets available work counters.  The supplied
\texttt{test\_c4\_cloning\_and\_startup.py} covers a minimal external
realization and preservation of the calibrated selector's type, record and
thresholds.  Unsupported external backends without the hook are rejected;
built-in backends have explicit cloning branches.  Fixed-step execution uses
the bound action context directly.

\subsection{Release-facing API}

The release distinguishes public availability from structural stability.  The
stable core includes problem, solution and state objects; \texttt{MethodSpec}
and catalogue accessors; fixed, generic adaptive and ETD34 solvers; startup
policies; basic analytic-expression constructors; the principal action
backends; calibrated selection; and workspace estimation.  External backends
enter the adaptive solvers through the \texttt{with\_tolerance} hook described
in Section~\ref{sec:adaptivity}.  Internal expression tree classes, the
backend-statistics dictionary beyond the \texttt{matvecs} entry, and selector
serialization remain public for research convenience but are not structurally
frozen.  Variable-step multistep and general linear methods and distributed
interfaces are explicitly outside the compatibility contract.

This release boundary is relevant to reproducibility: benchmark scripts may
inspect richer statistics, but the paper does not equate those internal keys
with a permanent user API.
